\documentclass[10pt,a4paper]{amsart}
\usepackage[margin=19mm]{geometry}
\usepackage{amsmath,amssymb,mathtools}
\usepackage[scr=rsfs,cal=euler]{mathalfa}
\usepackage{xfrac}
\usepackage[unicode,hidelinks,bookmarks=false]{hyperref}
\newcommand{\ie}{i.e.\ }
\newcommand{\cf}{cf.\ }
\newcommand{\resp}{respectively\ }
\newcommand{\Subsection}{\subsection}
\providecommand{\nobreakdash}{\nobreak\hbox{-}}
\setlength{\emergencystretch}{2em}
\multlinegap0pt
\usepackage{bm}
\usepackage{bbm}
\usepackage{dsfont}
\usepackage{xspace}
\usepackage{cancel}
\usepackage{mathtools}
\usepackage{amsthm}
\makeatletter
\@ifundefined{theorem}{\newtheorem{theorem}{Theorem}}{}
\makeatother
\def\b{\mathbf b}
\def\n{\mathbf n}
\def\rr{\mathbb L}
\def\r{\mathbb R}
\def\t{\mathbf t}
\title[Maximal translation surfaces in Lorentz-Minkowski space]{Maximal translation surfaces in Lorentz-Minkowski space}
\author{Rafael L\'opez}
\date{Source: arXiv:2507.14103v1 (CC BY 4.0)}
\begin{document}
\maketitle

\noindent\emph{Source and licence.} Maximal translation surfaces in Lorentz-Minkowski space. Version arXiv:2507.14103v1 (CC BY 4.0), distributed under the Creative Commons Attribution 4.0 International licence, as recorded in the arXiv licence metadata for this exact version and shown on the abstract page https://arxiv.org/abs/2507.14103v1. Full rights notice: licenses/arxiv-2507.14103-CC-BY-4.0.txt.\par\smallskip

\noindent\textit{Editorial scope.} Counted unit: the Theorem whose source label is \texttt{t31} in the article (source0.tex lines 265--278, source label \texttt{t31}), delivered together with the article's own complete original proof of it (source0.tex lines 280--343, one \texttt{proof} environment of 64 lines). No mathematics is written, repaired or re-derived: statement and proof are the article's own text line for line. Required context retained uncounted: tt (lines 95--97); t1 (lines 115--117); f1 (lines 164--170); f3 (lines 173--179); pr11 (lines 187--192); eq1 (lines 222--224). Every definition, display and label of those lines is retained unchanged. Omitted from this package and not redistributed: 1--94, 98--114, 118--163, 171--172, 180--186, 193--221, 225--264, 279--279, 344--878. Nothing inside a retained excerpt is omitted or altered: every retained body is delivered exactly as the article's lines are written, so no omission receipt of any kind accompanies this package. Citations: the retained excerpts carry no citation command at all, so no bibliography entry of the article is transcribed and no reference is redistributed. Reference adaptation: none was needed and none was made; every reference inside the delivered unit resolves inside it, so no unresolved reference or citation is produced. Printed slips kept literal: no printed word, symbol, hypothesis or number of the counted statement or of its proof was altered. Layout and class: the article's own class is \texttt{amsart} with packages including amssymb, amsthm, enumerate, diagbox, amscd, graphicx, color, url, relsize; the wrapper is the lane's pinned \texttt{amsart} editorial wrapper at 10pt on A4, which restates the theorem-like environments the retained text needs and adds the article's own macro definitions that the retained text needs (\texttt{\textbackslash b}, \texttt{\textbackslash n}, \texttt{\textbackslash r}, \texttt{\textbackslash rr}, \texttt{\textbackslash t}), with extra wrapper packages (bm, bbm, dsfont, xspace, cancel, mathtools). The delivered numbering is the unit's own. Assets: no retained span contains a figure environment, a graphics-inclusion command, a PDF-page inclusion, a table or a TikZ picture; no image file is copied into this package, the package declares no asset dependency and no image file is redistributed with it. Licence: version arXiv:2507.14103v1 of this article is distributed under the Creative Commons Attribution 4.0 International licence, as recorded in the arXiv licence metadata for this exact version and shown on the abstract page https://arxiv.org/abs/2507.14103v1. A full rights notice is delivered at licenses/arxiv-2507.14103-CC-BY-4.0.txt. Characters: every retained excerpt is pure ASCII, so the delivered engine, XeLaTeX with its default Latin Modern text font, reports no missing character.
\begin{equation}\label{tt}
X(s,t)=\alpha(s)+\beta(t).
\end{equation}

\begin{theorem} \label{t1}
Let $M$ be a maximal translation  surface in $\rr^3$. If a generating curve is planar, then the other generating curve is also planar.  
\end{theorem}

\begin{equation}\label{f1}
\left\{\begin{array}{ccccc}
\t'&=& &\kappa \n & \\
\n'&=&-\epsilon\kappa \t& +& \tau \b\\
\b'&=& & \tau \n. & \\
\end{array}\right.
\end{equation}

\begin{equation}\label{f3}
\left\{\begin{array}{ccccc}
\t'&=& & \n & \\
\n'&=&& \kappa \n&  \\
\b'&=&-\t &  &-\kappa \b. \\
\end{array}\right.
\end{equation}

\begin{equation}\label{pr11}
\begin{split}
\alpha(s)&=\frac{s^2}{2}\vec{v}+s\vec{b},\quad (\kappa=0),\\
\alpha(s)&=e^{ks}\vec{v}-\frac{s}{k}\vec{b},\quad (\kappa\not=0),
\end{split}
\end{equation}

\begin{equation}\label{eq1}
 (\t_\alpha,\t_\beta,\t_\alpha')+(\t_\alpha,\t_\beta,\t_\beta')=0.
\end{equation}

\begin{theorem} \label{t31}
 Let $M$ be a spacelike  translation surface \eqref{tt} such that  $\alpha''$  is spacelike and  $\beta$ is a pseudo-null curve. If $M$ is a maximal surface, then $M$ is a plane or both curves are planar curves. In the latter case,   after a rigid motion, the surface can be parametrized by 
\begin{equation}\label{e11}
X(x,t)=(x,f(x),0)+\beta(t),
\end{equation}
where  
\begin{equation}\label{e12}
 \begin{split}
 f'(x)&=-\tan\left(k(\cos\theta f(x)-x\sin\theta)+a\right), \\ 
 \beta(t)&=me^{kt}(-\sin\theta,\cos\theta,1)+t(\cos\theta,\sin\theta,0),
 \end{split}
 \end{equation}
 and $a,k,m,\theta\in\r$, $k,m\not=0$.
   \end{theorem}

\begin{proof}
Recall that $\beta$ is included in a plane because $\beta$ is a pseudo-null curve. Frenet equations \eqref{f1} and the identity  $\t_\beta\times\n_\beta=-\n_\beta$ imply that Eq. \eqref{eq1} becomes
\begin{equation}\label{c1}\kappa_\alpha\langle \b_\alpha,\t_\beta\rangle-\langle\t_\alpha,\n_\beta\rangle=0.
\end{equation}
The principal curvatures of $M$ are $\kappa_1=\kappa_\alpha\langle \b_\alpha,\t_\beta\rangle$ and $\kappa_2=-\langle\t_\alpha,\n_\beta\rangle$. Differentiating with respect to  $t$, we have
$$\kappa_\alpha\langle\b_\alpha,\n_\beta\rangle-\kappa_\beta\langle\t_\alpha,\n_\beta\rangle=0.$$
From   \eqref{c1}, we have 
\begin{equation}\label{c11}
\kappa_\beta\langle\b_\alpha,\t_\beta\rangle-\langle\b_\alpha,\n_\beta\rangle=0.
\end{equation}
Differentiating with respect to $t$ again, the equations of Frenet \eqref{f3} give
\begin{equation}\label{k21}\kappa_\beta'\langle\b_\alpha,\t_\beta\rangle=0.
\end{equation}
If $\langle\b_\alpha,\t_\beta\rangle=0$ identically, then the principal curvature $\kappa_1$ of $M$ is identically $0$ and   thus $M$ is a plane. 

Suppose that  $\langle\b_\alpha,\t_\beta\rangle\not=0$ at some point $(s_0,t_0)\in I\times J$. Then \eqref{k21} implies that  in an interval around $s_0$, the curvature $\kappa_\beta$ is constant.   We distinguish if $k$ is $0$ or not.
 
\begin{enumerate}
\item Case $k=0$. By  \eqref{pr11} we know that 
$$\beta(t)=\frac{t^2}{2}\vec{v}+\vec{b}t,$$
where $\vec{v},\vec{b}\in\rr^3$,  $\vec{v}$ is a spacelike vector and $\vec{b}$ is a unit spacelike vector orthogonal to $\vec{v}$.   Putting in \eqref{c1}, we have 
\begin{equation}\label{c8}
\kappa_\alpha\langle\b_\alpha,\vec{v}\rangle t+\kappa_\alpha\langle\b_\alpha,\vec{b}\rangle-\langle\t_\alpha,\vec{v}\rangle=0.
\end{equation}
This is a polynomial on $t$, hence $\langle\b_\alpha,\vec{v}\rangle=0$. Thus $\vec{v}$ is a linear combination of $\t_\alpha$ and $\n_\alpha$. Since $\vec{v}$ is lightlike and $\{\t_\alpha,\n_\alpha\}$ are orthonormal spacelike vectors, this combination is trivial, so $\vec{v}=0$. This gives a contradiction.
\item Case $k\not=0$. From the proof of   \eqref{pr11}, we know that 
\begin{equation}\label{c12}
\beta(t)=e^{kt}\vec{v}-\frac{t}{k}\vec{b},
\end{equation}
with $\vec{v}$  is lightlike orthogonal to $\vec{b}$ and $\vec{b}$ is a spacelike vector with $|\vec{b}|=k$. 
Substituting in \eqref{c1},  we have 
$$k e^{kt}\left(\kappa_\alpha \b_\alpha -k \t_\alpha,\vec{v}\rangle\right)-\frac{\kappa_\alpha}{k}\langle\b_\alpha,\vec{b}\rangle=0.$$
Therefore
\begin{equation}\label{c9}
\begin{split}
0&= \langle \kappa_\alpha \b_\alpha-k\t_\alpha,\vec{v}\rangle,\\
0&=\frac{\kappa_\alpha}{k}\langle\b_\alpha,\vec{b}\rangle.
\end{split}
\end{equation}
From the second equation, we have $\langle\b_\alpha,\vec{b}\rangle=0$. Differentiating, we get  $\tau_\alpha\langle\n_\alpha,\vec{b}\rangle=0$. We distinguish two cases.
\begin{enumerate}
\item Case $\langle\n_\alpha,\vec{b}\rangle=0$ identically. Differentiating again, we have $\kappa_\alpha\langle\t_\alpha,\vec{b}\rangle=0$. Since $\kappa_\alpha\not=0$, then $\langle \t_\alpha,\vec{b}\rangle=0$. This gives $\vec{b}=0$, a contradiction. 
\item Case $\tau_\alpha=0$.   In particular, $\alpha$ is a planar curve. This proves Thm. \ref{t1} in this situation. After a rigid motion, we can assume that $\alpha$ is contained in the $xy$-plane and  $\b_\alpha=-(0,0,1)$. Then there are $m,\theta\in\r$ such that
\begin{equation*}
\begin{split}
\vec{v}&=m(-\sin\theta,\cos\theta,1),\\
 \vec{b}&=k(\cos\theta,\sin\theta,0).
 \end{split}
 \end{equation*}
Both vectors in \eqref{c12} provide the parametrization of $\beta$  in \eqref{c9}. We now parametrize $\alpha$ by $\alpha(x)=(x,f(x),0)$ for some smooth function $f=f(x)$. Then 
$$\kappa_\alpha=\frac{f''}{(1+f'^2)^{3/2}},\quad \t_\alpha=\frac{(1,f',0)}{\sqrt{1+f'^2}}.$$ 
Then
$$\kappa_\alpha\b_\alpha-k\t_\alpha=\left(-\frac{k}{\sqrt{1+f'^2}},\frac{kf'}{\sqrt{1+f'^2}},-\frac{f''}{(1+f'^2)^{3/2}}\right).$$
The first equation of \eqref{c9} gives 
$$\frac{f''}{1+f'^2}=k\sin\theta+kf'\cos\theta.$$
A first integration of this equation yields \eqref{e12}.   


 
\end{enumerate}


\end{enumerate}
\end{proof}

\end{document}
